\begin{lemma}[Situation 7.2, Subcase 1.1 stored exception]
Let \(D\) be a \noderef{HierarchySituationTwoCaseOneData}.  Suppose
there is a color \(\eta_h\in D_{n,q}\) such that
\[
       \theta\in\Gamma^q_{\eta_h}
       \quad\hbox{and}\quad
       \eta_h\in\overline\varphi_n(y_{p-1}).
\]
Then a contradiction follows.
\end{lemma}
\begin{proof}
Put \(H=C.\mathit{hierarchy}\), let \(q\) be its last divider index,
put \(A=\noderef{HierarchyStarredStage}(H,q)\), and set
\[
 \ell=\noderef{HierarchyFullTailLength}(C),\quad
 B=\noderef{HierarchyFullTailPrefix}(C,\ell-2),\quad
 B^+=\noderef{HierarchyFullTailPrefix}(C,\ell-1).
\]
Apply \noderef{HierarchyCounterexamplePenultimateElementary} to the
six relevant fields of
\noderef{HierarchyMinimumCounterexampleContext}.  Thus
\[
                  V(B^+)\text{ is elementary under }\varphi_n.       \tag{1}
\]
Because \(\alpha,\eta_h\in D_{n,q}\) and both are missing at
\(y_{p-1}\), (1) and the exact chain \(A\preccurlyeq B\preccurlyeq
B^+\) make \(y_{p-1}\) their first missing vertex.  The avoidance
field of \noderef{GoodHierarchy}, applied to
\(\Gamma^q_\alpha=\{\gamma_1,\gamma_2\}\) and to
\(\theta\in\Gamma^q_{\eta_h}\), gives
\[
 \begin{aligned}
 &\gamma_1,\gamma_2\notin
       \varphi_n\langle B^+-A\rangle,\\
 &\theta\notin\varphi_n\langle B^+-A\rangle.                 \tag{2}
 \end{aligned}
\]
For \(\alpha\), the same tail exclusion holds above \(A\), except for
the exact divider-zero SE/RE seed when it carries the stored color
\(\alpha\).  In that exception the singleton connecting-edge equality
identifies the one protected edge; no other parallel edge is exempted.

Apply \noderef{ChangeLemma} with the ordered pair
\((\gamma_j,\alpha)\), using (1), (2), the stage-missing half of
avoidance, and the minimum-context fields.  For \(j=1,2\) it gives
\[
 K_{\alpha,\gamma_j}(v_{\gamma_j})
      =K_{\alpha,\gamma_j}(y_{p-1}),                          \tag{3}
\]
where \(v_{\gamma_j}\in V(A)\) is the first
\(\gamma_j\)-missing vertex.

The two members of \(\Gamma^q_\alpha\) are distinct.  At least one is
different from \(\theta\); rename the displayed pair, if necessary, so
that
\[
                           \theta\ne\gamma_2.                 \tag{4}
\]
Also \(\theta\ne\alpha\): the exact edge \(e_p\) is incident with
\(y_p\) and has color \(\theta\), whereas \(\alpha\) is missing there.
The component in (3) for \(j=2\) is disjoint from
\(Q_1=K_{\alpha,\gamma_2}(y_p)\).  Indeed, \(\alpha\) is missing at
both \(y_{p-1}\) and \(y_p\), while \(\gamma_2\) is missing at
\(v_{\gamma_2}\).  These three vertices are distinct by the vertex
injectivity field of \noderef{TreeSequence}.  If the components were
equal, \noderef{TwoColorComponentSpanningPath} and
\noderef{TwoColorIncidenceDegreeAtMostTwo} would put three
degree-at-most-one vertices on one path.  Hence
\noderef{TwoColorComponentsDisjoint} gives the asserted disjointness,
and \noderef{TwoColorComponentSpanningPath} supplies the required
\noderef{TwoColorPathComponent} certificate rooted at \(y_p\), including
the possible isolated-vertex case.

Define the first recoloring explicitly by
\[
 \mu_1=\noderef{TwoColorComponentSwap}
        (G,\varphi_n,\gamma_2,\alpha,\operatorname{inl}y_p).  \tag{5}
\]
Before using the full-tail provider, apply
\noderef{StableChangeCasesOneTwo} with the ordered colors
\((\gamma_2,\alpha)\), (1), (3), the path certificate for \(Q_1\), and
the Cases 1/2 alternative supplied by the Gamma stage-missing property:
when \(q>0\), \(\gamma_2\) is missing on
\(A=T_{n,q}\), while when \(q=0\) the zero-index avoidance clause puts
\(\gamma_2\) among the colors missing on \(T_n\).  Retain in particular its
\noderef{HierarchyWeaklyStable} conclusion for \(\mu_1\).

We separately prove missing-set equality on \(B^+\).  Fix
\(u\in V(B^+)\).  If its vertex-node is outside \(Q_1\), the swap changes
no edge incident with \(u\).  If it belongs to \(Q_1\), it cannot be an
endpoint of the spanning path: an endpoint vertex-node has incidence
degree at most one, so one of \(\alpha,\gamma_2\) is missing there; by
elementarity (1), its component would then be the component in (3),
contrary to disjointness.  Hence it is an internal degree-two vertex.
Unfolding \noderef{TwoColorIncidenceGraph}, properness says that its two
incident edge-nodes are exactly one \(\alpha\)-edge and one
\(\gamma_2\)-edge.  Swapping them preserves the set of present colors,
and hence the missing-color set, at \(u\).

We also prove displayed-edge equality independently.  On \(B^+-A\),
(2) excludes \(\gamma_2\).  Apply
\noderef{HierarchyTailMissingColorExclusion} to \(\alpha\), which is
missing at \(y_{p-1}\).  It excludes \(\alpha\) on this interval except
possibly when \(q=0\), \(A=T_n\), the actual terminal
\noderef{TashkinovSeriesStep} is SE or RE, and the edge is the individual
singleton seed \(f_n\) with \(\varphi_n(f_n)=\alpha\in D_n\).
In precisely this exception, unfold
\noderef{HierarchyWeaklyStable}: since the kind is non-PE it is
\((T_n,D_n,\varphi_n)\)-\noderef{StableColoring}.  The exact adjoin
certificate makes \(f_n\) incident with \(T_n\) in the sense of
\noderef{TreeIncidentEdge}; the edge clause of stable coloring therefore
fixes this individual seed.  Finally \(e_p\) has color \(\theta\),
distinct from both swapped colors by (4) and by properness at \(y_p\).
Therefore
\[
 \begin{aligned}
 &\mu_1(T.\mathit{edge}(r))=\varphi_n(T.\mathit{edge}(r))
       &&(|A|\leq r<|T|),\\
 &\overline\mu_1(u)=\overline\varphi_n(u)
       &&(u\in V(B^+)).
 \end{aligned}                                                \tag{6}
\]

The totalized first-missing position of every stored color is unchanged.
Indeed, unfold \noderef{HierarchyStoredColorsAt} and
\noderef{HierarchyFirstMissingPosition}.  The color \(\gamma_2\) is not
stored because it is missing on \(A\).  Every stored color other than
\(\alpha\) is unaffected at every vertex.  The color \(\alpha\) has no
missing witness on \(A\), while (6) identifies its complete witness set
through \(B^+\) and that set contains \(y_{p-1}\); hence the two least
positions are equal.  Apply
\noderef{StableChangeCasesOneTwoFullTail}, ordered as
\((\gamma_2,\alpha)\), with the exact identities
(2), (3), and (6).  Its conclusion supplies weak stability,
\noderef{ModColoring}, an actual corresponding ETT on \(T\), maximum
property, a rebuilt base and full \noderef{GoodHierarchy}, and
\noderef{SameGoodHierarchy}.  Compose the returned correspondence with
the identity presentation of the original ETT, and use the transitivity
of stable colorings from \noderef{sc2}; copying the context fields gives
a genuine \noderef{HierarchyReplacementPresentation} \(P_1\) whose
coloring is \(\mu_1\).  More explicitly, the returned objects fill its
correspondence, maximum, base, hierarchy, goodness, and same-good fields;
\(v_{\gamma_2}\) and \(y_p\) are two distinct missing witnesses and fill
the non-elementarity field.  Properness is supplied by
\noderef{TwoColorComponentSwapProper}, and the swap preserves the exact
uncolored-edge set, so \noderef{IsOneEdgeUncolored} is retained.  Copy
the rung-uniform lower assertion and fewer-divider fields.  Same-good
identifies the starred stage and divider count, while the terminal tree
is literally \(T\), so the counterexample measures are equal before the
old lexicographic-minimality field is transported.  Finally weak
stability gives equality of missing sets on \(T_n\), transporting the
large-set bound.  These are all fields of the replacement
\noderef{HierarchyMinimumCounterexampleContext}, rather than an
unjustified wholesale copy.  In particular (6) transports (1), and
\(\gamma_2\) is missing at both \(v_{\gamma_2}\) and \(y_p\).

We now make the manuscript's counting choice.  Apply \noderef{NineN}
to \(P_1\) at \(i=\ell-2\).  Its available set consists of colors
missing on \(B\), with the base-extra missing colors and all colors used
on the exact prefix \(B-A\) removed, and it has cardinality at least
\(2n+11\).  Passing from \(B\) to \(T\) adds exactly the individual
edges \(e_{p-1},e_p\), so excluding every color used on \(T-A\) removes
at most two more colors.  By
\noderef{HierarchyStoredColorsCardinality},
\[
 |\Gamma^q|\leq2|D_{n,q}|\leq2n.                            \tag{7}
\]
The resulting set is nonempty.  Choose
\[
 \beta\in\overline\mu_1(B)
   \setminus\overline\mu_1(\mathcal B_1.\mathit{core}-V(T_n^1))
   \setminus\mu_1\langle T-A\rangle
   \setminus\Gamma^q.                                       \tag{8}
\]
This records the selected color itself, not merely its cardinality
bound.  In particular \(\beta\) is missing first at a vertex
\(v_\beta\in V(B)\), it is not a Gamma color, and neither \(\beta\) nor
\(\gamma_2\) is used on the exact tail segment relevant below.

Apply \noderef{ChangeLemma} in the presentation \(P_1\), with the
ordered pair \((\gamma_2,\beta)\).  Equations (6)--(8), the copied minimum context,
and the copied Gamma data give
\[
 K_{\beta,\gamma_2}(v_\beta)
       =K_{\beta,\gamma_2}(v_{\gamma_2}).                    \tag{9}
\]
This component is disjoint from
\(Q_2=K_{\beta,\gamma_2}(y_p)\), as follows.  Split explicitly on
\(v_\beta=v_{\gamma_2}\).  In the equality branch both \(\beta\) and
\(\gamma_2\) are missing at the common vertex.  Unfolding
\noderef{MissingColors} and \noderef{TwoColorIncidenceGraph}, its
vertex-node has no incident edge-node, so its component is the isolated
singleton and cannot contain the distinct vertex \(y_p\).  In the
inequality branch \(v_\beta,v_{\gamma_2},y_p\) are three distinct
vertices of incidence degree at most one, because respectively
\(\beta,\gamma_2,\gamma_2\) is missing there.  If their components were
equal, \noderef{TwoColorComponentSpanningPath} would put three endpoints
on a single spanning path, contrary to
\noderef{TwoColorIncidenceDegreeAtMostTwo}.  Thus
\noderef{TwoColorComponentsDisjoint} gives disjointness in both branches,
and the spanning-path theorem supplies the possibly singleton
path-component certificate for \(Q_2\).
Define the second
recoloring
\[
\mu_2=\noderef{TwoColorComponentSwap}
        (G,\mu_1,\gamma_2,\beta,\operatorname{inl}y_p).       \tag{10}
\]
Apply \noderef{StableChangeCasesOneTwo} first, relative to \(P_1\), and
use the same explicit alternative: \(\gamma_2\) is missing on
\(T_{n,q}\) if \(q>0\), and on \(T_n\) if \(q=0\).  Retain its
weak-stability conclusion.  The endpoint/internal-vertex
argument used for (6), now using (9), proves missing-set equality on
\(B^+\).  Equation (8), the avoidance of \(\gamma_2\), and the
individual-seed edge clause of that weak stability prove that every exact
displayed edge at a position from \(|A|\) through \(|T|-1\) is fixed.

For the stored-position premise, unfold
\noderef{HierarchyStoredColorsAt} and
\noderef{HierarchyFirstMissingPosition}.  The color \(\gamma_2\) is not
stored because it is missing on \(A\).  If \(\beta\) is not stored, all
stored colors are unaffected.  If it is stored, it has no witness on
\(A\), while (8) gives a witness in \(B\); missing-set equality on
\(B^+\) identifies the complete candidate sets up to that witness, and
therefore identifies their least positions.  Hence every stored
first position is unchanged.  A second application of
\noderef{StableChangeCasesOneTwoFullTail}, ordered as
\((\gamma_2,\beta)\) and now relative to \(P_1\),
gives an actual corresponding MP ETT, rebuilt full good hierarchy, and
same-good data under \(\mu_2\).  Composing both correspondence
certificates and both same-good relations, and using \noderef{sc2} at
each positive rung, produces an original-relative replacement
presentation \(P_2\).  Field by field, its non-elementarity uses the
distinct witnesses \(v_\beta,y_p\); properness and the uncolored-edge
field use \noderef{TwoColorComponentSwapProper} and the definition of
the swap; the lower-rung and fewer-divider fields copy because the rung
is unchanged; literal equality of \(T\) and same-good give equality of
the two measures before minimality is transported; and weak stability
transports the large-set bound through missing-set equality on \(T_n\).
It thus carries the full transported minimum context, and \(\beta\) is
missing at both \(v_\beta\) and \(y_p\).

The selected stored color \(\eta_h\) is still missing at
\(y_{p-1}\), since (6) and its analogue for (10) preserve all missing
sets on \(B^+\).  Its Gamma pair is unchanged and contains \(\theta\).
The avoidance field therefore gives
\[
 \theta\notin\mu_2\langle B^+-A\rangle.                      \tag{11}
\]
Equation (8) gives the stronger
\[
 \beta\notin\mu_2\langle T-A\rangle.                         \tag{12}
\]
Apply \noderef{ChangeLemma} under \(P_2\), with the ordered pair
\((\theta,\beta)\).  The first \(\theta\)-missing vertex lies in the
unchanged last ordinary stage, while \(v_\beta\in B\); (11)--(12) give
the exact unused intervals.  Split on \(v_\beta=v_\theta\).  In the
equality branch both colors are missing at the common vertex, so
unfolding \noderef{MissingColors} and
\noderef{TwoColorIncidenceGraph} shows that its vertex-node has no
neighbor and its component is isolated.  In the inequality branch
\(v_\beta,v_\theta,y_p\) are three distinct degree-at-most-one vertices,
because \(\beta,\theta,\beta\) is missing at them respectively; a common
component would contradict the spanning path of
\noderef{TwoColorComponentSpanningPath}.  Hence in either branch their
component is disjoint from \(Q_3=K_{\beta,\theta}(y_p)\), and the same
spanning-path theorem supplies its possibly singleton path-component
certificate.  Define the third recoloring
\[
 \mu_3=\noderef{TwoColorComponentSwap}
        (G,\mu_2,\theta,\beta,\operatorname{inl}y_p).         \tag{13}
\]
Apply \noderef{StableChangeCasesOneTwo} to the presentation \(P_2\),
with \(\theta\) as the stage-missing affected color and \(\beta\) as
the other color.  The Gamma avoidance field puts \(\theta\) among the
colors missing on \(T_{n,q}\) when \(q>0\), and on \(T_n\) when \(q=0\),
which is exactly its Cases 1/2 alternative.  The component separation just proved is its
disjointness premise; (11)--(12) give its exact unused interval; the
properness, one-uncoloredness, lower-rung, fewer-divider, (6.6),
large-set, and minimality premises are the fields of \(P_2\)'s minimum
context.  Retain in particular its
\noderef{HierarchyWeaklyStable} conclusion for \(\mu_3\).

We next extract the actual terminal construction of the replacement
presentation.  Put
\[
 \mathcal E_2=P_2.\mathit{corresponding}.\mathit{replacementETT},
 \qquad
 \mathcal S_2=P_2.\mathit{corresponding}.\mathit{replacementSeries},
\]
and let \(L_2\) be the last output tree of this series and \(T_n^2\) its
terminal ladder tree.  The \(\mathit{seriesPresentation}\) field of
\noderef{CorrespondingExtendedTashkinovTree} transports the dependent
series of \(\mathcal E_2\) to \(\mathcal S_2\), and its
\(\mathit{finalGeneratingColoring}\) field identifies the coloring at the
terminal series index with \(\mu_2\), after casting that index along the
\(\mathit{rungAgreement}\) field.  Thus the certificate
\(\mathcal E_2.\mathit{series}.\mathit{certified}\) in the actual
\noderef{TashkinovSeries} may be inspected at its last index.  The
initial constructor is impossible by the positive-rung field of
\noderef{ExtendedTashkinovTree}; unpacking the remaining actual
\noderef{TashkinovSeriesStep} gives the following exhaustive alternatives.

In the PE alternative, its closed-extension field is a
\(\mu_2\)-\noderef{TashkinovClosure} of \(T_n^2\) to \(L_2\).  Unfolding
that closure gives a
\(\mu_2\)-\noderef{TashkinovAugmentationSequence}
\(T_n^2\longrightarrow L_2\).  In the SE and RE alternatives, the
constructor supplies an individual edge \(f_n^2\), the literal singleton
connecting-edge equality, a tree \(S_2\) with
\noderef{TreeSequenceAdjoinEdge}\((T_n^2,f_n^2,S_2)\), and a
\(\mu_2\)-\noderef{TashkinovClosure} from \(S_2\) to \(L_2\); hence it
supplies a \(\mu_2\)-augmentation sequence
\(S_2\longrightarrow L_2\).  The extension-kind and connecting-edge
fields of the same correspondence transport these exact alternatives
back to the original-relative terminal arrays, while its
\(\mathit{ladderTreesAgree}\) field identifies \(T_n^2\) with the
original-relative terminal ladder tree.  In what follows write
\(S\) for \(T_n^2\) in the PE alternative and for this exact seed
\(S_2\) in the SE/RE alternatives, and write
\(\rho:S\longrightarrow L_2\) for the resulting certified route.

We now prove, case by case, the first exact prefix needed to restrict
\(\rho\).  Transport the stages, core, and divider count through
\(P_2.\mathit{sameGood}\), and call the replacement hierarchy \(H_2\).
Thus its last starred stage is the same exact tree \(A\).
If the terminal kind is PE and \(q>0\), the
\(\mathit{baseStage}\) equality and the iterated
\(\mathit{nested}\) fields of \noderef{TashkinovHierarchy} give
\[
                      T_n^2\preccurlyeq A\preccurlyeq B^+.
\]
If the kind is PE and \(q=0\), the
\(\mathit{parallel}\) field of the replacement
\noderef{HierarchyBase} supplies its closure \(R_n^2\) and
\[
 U_2:\noderef{OrderedTreeUnion}(T_n^2,R_n^2,A).
\]
The field \(U_2.\mathit{leftSegment}\), followed by
\(A\preccurlyeq B^+\), again gives \(S=T_n^2\preccurlyeq B^+\).

Suppose next that the kind is SE or RE and \(q=0\).  The nonparallel
base clause gives \(A=T_n^2\), while the exact adjoin certificate gives
\(|S|=|A|+1\) and \(A\preccurlyeq S\).  The hypothesis
\(D.\mathit{tailAtLeastTwo}\), together with the definitions of
\(\ell\) and \(B^+\), gives \(|A|+1\leq|B^+|\).  The closure route makes
\(S\) an exact prefix of \(L_2\): induct on its
\noderef{TashkinovAugmentationSequence} and compose the exact-prefix
clause of each \noderef{TashkinovAugmentation}.  Also
\(B^+\preccurlyeq T\preccurlyeq L_2\).  Comparing the entries of these
two prefixes of \(L_2\), using the displayed length inequality and
\noderef{TreeSequenceInitialSegment}, proves \(S\preccurlyeq B^+\).

Finally suppose that the kind is SE or RE and \(q>0\).  The base-stage
identity gives \(H_2.\mathit{stage}(0)=T_n^2\), and the first strict
nesting gives
\[
 |H_2.\mathit{stage}(1)|\geq |T_n^2|+1=|S|.
\]
Both \(S\) and \(H_2.\mathit{stage}(1)\) are exact prefixes of the same
tree \(L_2\): for \(S\) this follows by induction along its closure
route, and for stage one it follows from the iterated hierarchy nesting,
the terminal-stage equality, and
\(\mathcal E_2.\mathit{belowFinal}\).  Moreover
\noderef{HierarchyNonparallelFirstEdge}, applied to \(H_2\) and the
literal singleton edge equality, identifies the entry of stage one at
position \(|T_n^2|\) with the same individual edge \(f_n^2\) adjoined in
\(S\); no parallel edge is substituted.  Comparing all earlier entries
in the common final output now gives
\[
 S\preccurlyeq H_2.\mathit{stage}(1)
       \preccurlyeq A\preccurlyeq B^+.
\]
This completes the proof of \(S\preccurlyeq B^+\) in every terminal-kind
and divider case.

The other exact prefix required for restriction is independent of these
cases.  The definition of \(B^+\) gives \(B^+\preccurlyeq T\), and the
\(\mathit{belowFinal}\) field of
\noderef{ExtendedTashkinovTree} for \(\mathcal E_2\) gives
\(T\preccurlyeq L_2\).  Composing the length inequalities and the
entrywise vertex and edge identities gives
\[
                         B^+\preccurlyeq L_2.                  \tag{14a}
\]
Apply \noderef{TashkinovAugmentationSequenceRestrict} to \(\rho\) with
the two now proved premises \(S\preccurlyeq B^+\) and (14a).  This gives
the old-color route \(S\longrightarrow B^+\).

Separately, compose \(S\preccurlyeq B^+\preccurlyeq T\) and use again
\(T\preccurlyeq L_2\).  A second application of
\noderef{TashkinovAugmentationSequenceRestrict}, now with the exact
intermediate tree \(T\), gives a \(\mu_2\)-route
\(S\longrightarrow T\).  We obtain its suffix from \(B^+\) directly,
not from the restriction theorem.  Unfold
\noderef{TashkinovAugmentationSequence} as
\(\operatorname{Relation.ReflTransGen}\) and induct on its constructors.
Every one-step \noderef{TashkinovAugmentation} raises the length by
exactly one and makes its predecessor an exact prefix of its successor;
induction also makes every intermediate tree an exact prefix of the
common endpoint \(T\).  Compare the length of each successive member
with \(|B^+|\).  Since \(|S|\leq|B^+|\), while the definitions of the
full-tail prefixes and \(D.\mathit{tailAtLeastTwo}\) give
\(|T|=|B^+|+1\), a first member of length at least \(|B^+|\) exists.
The unit length increase makes its length exactly \(|B^+|\).  That member
and \(B^+\) are equal-length exact prefixes of \(T\), so their vertex and
edge entries agree at every index by
\noderef{TreeSequenceInitialSegment}, and proof irrelevance identifies
the tree-sequences.  Splitting the reflexive-transitive chain at this
member yields the required old-color suffix
\[
                         B^+\longrightarrow T.                \tag{14b}
\]

It remains to transport only the prefix route.  First, missing sets agree
on every vertex of \(B^+\).  Indeed, reapplying
\noderef{HierarchyCounterexamplePenultimateElementary} to
the six relevant fields of \(P_2.\mathit{minimumContext}\), and using
\(P_2.\mathit{sameGood}\) to identify the resulting full-tail prefix,
makes \(B^+\) elementary under \(\mu_2\).
If a vertex-node of \(B^+\) is outside \(Q_3\), no incident edge changes.
If it lies in \(Q_3\), it cannot be an endpoint: an endpoint is missing
one of \(\beta,\theta\), while elementarity and the retained missing
witness for that color in the component disjoint from \(Q_3\) force that
endpoint to be the retained witness.  It therefore cannot lie in
\(Q_3\).  Thus it is an internal degree-two vertex.  Properness gives
exactly one incident \(\beta\)-edge and one incident \(\theta\)-edge;
the swap exchanges their values and fixes every other color.  Hence the
set of present colors, and therefore the missing set, is unchanged.

For the fixed-edge premise, take an exact position \(t\) of \(B^+\) with
\(|S|\leq t\).  If \(|A|\leq t\), equations (11)--(12) exclude both
swapped colors on that displayed edge, so (13) fixes it.  Suppose
\(t<|A|\).  If \(q>0\), unfolding
\noderef{HierarchyWeaklyStable} gives the
\noderef{TreeColoringInvariant} clause on \(A\).  Since every tree source
has positive length by \noderef{TreeSequence}, \(|S|\leq t\) excludes
the root position, and that invariance fixes the displayed edge.  If
\(q=0\) and the kind is PE,
then \(S=T_n^2\); the inequalities
\(|T_n^2|\leq t<|A|\) and
\(U_2.\mathit{newEntryFromRight}\) identify the edge with an exact edge
of \(R_n^2\).  Unfolding weak stability gives the combined-stability
hypothesis of \noderef{CombinedStableClosurePreserved}.  The properness
and exact-one-uncoloredness assumptions for \(\mu_2\) are fields of
\(P_2.\mathit{minimumContext}\); those for \(\mu_3\) were established
from \noderef{TwoColorComponentSwapProper} and preservation of the
uncolored set before applying stable change.  The cited closure theorem
therefore fixes every displayed edge of \(T_n^2\) and \(R_n^2\), hence
this edge.  If \(q=0\)
and the kind is SE or RE, then \(A=T_n^2\) and the source is the exact
seed with \(|S|=|A|+1\), contradicting \(|S|\leq t<|A|\).  Notice in
particular that transport asks for no equality at the seed position
\(|A|\).  We have proved exactly
\[
 \begin{aligned}
 &\overline\mu_3(u)=\overline\mu_2(u)
       &&(u\in V(B^+)),\\
 &\mu_3(B^+.\mathit{edge}(t))=\mu_2(B^+.\mathit{edge}(t))
       &&(|S|\leq t<|B^+|).
 \end{aligned}                                                \tag{14}
\]
There is no claim that all edges of \(B^+\), or any later tail edge, are
fixed.

Apply \noderef{TashkinovAugmentationSequenceColoringTransport} to the
old route \(S\longrightarrow B^+\), with penultimate prefix \(B\).  The
exact prefix \(B\preccurlyeq B^+\) and equality
\(|B^+|=|B|+1\) follow from the full-tail-prefix definitions and
\(D.\mathit{tailAtLeastTwo}\); the first line of (14), restricted to
\(V(B)\), is its missing-set premise, and the second line of (14) is
exactly its edge-value premise.  We obtain the required actual
\(\mu_3\)-route \(S\longrightarrow B^+\).  Finally, the exact edge
\(e_p\) has value \(\theta\) under \(\mu_2\) and is incident with the
root \(y_p\) of \(Q_3\), so the component swap (13) gives
\[
                    \mu_3(e_p)=\beta.                         \tag{15}
\]
Both \(\beta\) and \(\theta\) remain missing on \(B^+\) under
\(\mu_3\), at their two component-disjoint witnesses.  Equation (14)
does not by itself settle the stored-position premise.  Unfold
\noderef{HierarchyStoredColorsAt} and
\noderef{HierarchyFirstMissingPosition}.  The color \(\theta\) is not
stored because it is missing on \(A\).  If \(\beta\) is stored, (8)
places its first witness in \(B\), and (14) identifies the candidate
sets through that witness, so the two least positions agree; if it is
not stored, all stored colors are unaffected.  For the Gamma cutoff, \(\beta\)
belongs to no Gamma pair by (8), while any pair containing \(\theta\)
is, by pairwise disjointness, the displayed pair for \(\eta_h\), whose
first missing vertex is \(y_{p-1}\in B^+\).  All hypotheses of
\noderef{StableChangeCasesOneTwoSplitTail} are therefore satisfied with
split \(X=B^+\).  It rebuilds the full exact tree \(T\) under \(\mu_3\),
with an actual corresponding ETT, maximum property, full good hierarchy,
and same-good data.  Compose these with \(P_2\) exactly as above to
obtain a replacement presentation \(P_3\) and its transported minimum
context.  Explicitly, the returned objects fill correspondence, maximum
property, base, hierarchy, goodness, and same-good; the retained
\(\theta\)-missing witnesses \(v_\theta,y_p\) fill non-elementarity.
\noderef{TwoColorComponentSwapProper} and preservation of the exact
uncolored set fill the first context field.  Compose correspondence and
same-good entry by entry, copy the rung-uniform lower and fewer-divider
fields, prove equality of measure from literal equality of \(T\) and
same-good before transporting minimality, and use weak stability on
\(T_n\) for the large-set bound.  Thus all six minimum-context fields
are present.

Form a new \noderef{HierarchySituationTwoTailData} \(D_3\) for \(P_3\).
Use the same structural indices and non-incidence equation.  Same-good
data retain the last starred stage, so the path number and full-tail
length remain those in \(D\).  Use \(\beta\) as its last-edge color;
(15) is the color equation.  It is present at \(y_{p-1}\), since
\(v_\beta\in V(B)\) is a distinct vertex missing \(\beta\) and (14)
transports elementarity of \(B^+\).  Thus every field of \(D_3\) is
verified.

Apply \noderef{HierarchySituationTwoReorderedTreesCertified} to
\(P_3\) and \(D_3\).  It supplies the actual reordered tree \(T^-_3\),
its corresponding same-rung ETT, maximum property, an actual replacement
\noderef{HierarchyBase}, and an actual \noderef{TashkinovHierarchy}
whose internal stages and divider positions are the old ones.  Construct
goodness pointwise.  Stored sets and Gamma pairs through every ordinary
stage copy.  For a stored color \(\xi\), if it has a missing witness in
\(B\), its least witness and all preceding displayed entries are
unchanged in \(T^-_3\).  If it has no witness in \(B\), its first
position in \(T^-_3=B+(e_p,y_p)\) is the final position, whether \(y_p\)
is a witness or the totalized fallback is used.  Deleting
\(e_{p-1}\) only removes avoidance constraints, and the only edge newly
placed after \(B\) is \(e_p\), whose color
\(\beta\notin\Gamma^q\).  Thus each Gamma color avoids the exact interval
beginning at its stored color's recalculated least position.  All other
good-hierarchy fields copy, so \(T^-_3\) has a full good hierarchy.

Finally, the swap (13) makes \(\theta\) missing at \(y_p\), while its
component-disjoint first witness in \(A\) remains missing.  Both
distinct vertices lie in \(T^-_3\).  Its vertex set is therefore not
elementary, and the corresponding ETT, maximum property, and copied
good hierarchy form a \noderef{HierarchyCounterexample} \(C^-_3\).
It has the same rung and divider count.  The reordered certificate gives
\(|T^-_3|=|T|-1\), while its last starred stage is still \(A\); hence its
full-tail length is exactly \(\ell-1\).  By the finite-minimum definition
of \noderef{HierarchyPathNumber}, its path number is at most its
full-tail length, so it is at most \(\ell-1<\ell\).  The datum \(D\)
says that the old path number is exactly \(\ell\).  Therefore the first
coordinate of \noderef{HierarchyCounterexampleMeasure} strictly
decreases, and
\noderef{HierarchyCounterexampleLess}(C^-_3,C)\), contrary to the
lexicographic-minimality field of the original context.  This is the
required contradiction.
\end{proof}
